
\documentclass{article}
\usepackage{colortbl}
\usepackage{makecell}
\usepackage{multirow}
\usepackage{supertabular}

\begin{document}

\newcounter{utterance}

\centering \large Interaction Transcript for game `cladder', experiment `full\_v1.5\_default', episode 8837 with selfrefine{-}Qwen3.5{-}9B.
\vspace{24pt}

{ \footnotesize  \setcounter{utterance}{1}
\setlength{\tabcolsep}{0pt}
\begin{supertabular}{c@{$\;$}|p{.15\linewidth}@{}p{.15\linewidth}p{.15\linewidth}p{.15\linewidth}p{.15\linewidth}p{.15\linewidth}}
    \# & \multicolumn{2}{c}{Player} && \multicolumn{2}{c}{Game Master} \\
    \hline

    \theutterance \stepcounter{utterance}  
    & & & \multicolumn{4}{p{0.6\linewidth}}{
        \cellcolor[rgb]{0.9,0.9,0.9}{
            \makecell[{{p{\linewidth}}}]{
                \texttt{\tiny{[P1$\langle$GM]}}
                \texttt{You are an expert in causal inference. The following question is not a typical commonsense query, but rather a meticulously designed question created by a professor specializing in causal inference, intended to assess the students' mastery of the course content.} \\
\texttt{Imagine a self{-}contained, hypothetical world with only the following conditions, and without any unmentioned factors or causal relationships: Yupt has a direct effect on zupj and muvq. Xyfo has a direct effect on zupj and muvq. Zupj has a direct effect on muvq. Xyfo is unobserved. For those who are not yupt and are not zupj, the probability of muvq is 44\%. For those who are not yupt and are zupj, the probability of muvq is 51\%. For those who are yupt and are not zupj, the probability of muvq is 57\%. For those who are yupt and are zupj, the probability of muvq is 10\%. For those who are not yupt and are not xyfo, the probability of zupj is 42\%. For those who are not yupt and are xyfo, the probability of zupj is 87\%. For those who are yupt and are not xyfo, the probability of zupj is 26\%. For those who are yupt and are xyfo, the probability of zupj is 44\%. The overall probability of xyfo is 22\%. Does yupt negatively affect muvq through zupj?} \\
\texttt{Start your answer with 'Yes' or 'No', followed by additional reasoning or evidence to support your explanation.} \\
            }
        }
    }
     \\ \\

    \theutterance \stepcounter{utterance}  
    & & \multicolumn{4}{p{0.6\linewidth}}{
        \cellcolor[rgb]{0.95,0.95,0.95}{
            \makecell[{{p{\linewidth}}}]{
                \texttt{\tiny{[P1$\leftrightarrow$P]}}
                \texttt{{[}DRAFT{]}} \\
\texttt{No} \\
\\ 
\texttt{**Reasoning:**} \\
\\ 
\texttt{To determine if \$Y\_\{upt\}\$ negatively affects \$M\_\{uvq\}\$ through \$Z\_\{upj\}\$, we must calculate the **Natural Direct Effect (NDE)** and the **Natural Indirect Effect (NIE)** of \$Y\_\{upt\}\$ on \$M\_\{uvq\}\$ via \$Z\_\{upj\}\$. Specifically, we need to compare the probability of \$M\_\{uvq\}\$ when \$Y\_\{upt\}\$ changes from 0 to 1, while holding the mediator \$Z\_\{upj\}\$ constant at the value it would have taken under the baseline condition (\$Y\_\{upt\}=0\$).} \\
\\ 
\texttt{The causal graph structure is:} \\
\texttt{\$X\_\{yfo\} \textbackslash{}to Z\_\{upj\} \textbackslash{}to M\_\{uvq\}\$} \\
\texttt{\$Y\_\{upt\} \textbackslash{}to Z\_\{upj\}\$} \\
\texttt{\$Y\_\{upt\} \textbackslash{}to M\_\{uvq\}\$} \\
\texttt{(\$X\_\{yfo\}\$ is unobserved, but we have sufficient data to compute the necessary conditional probabilities).} \\
\\ 
\texttt{**Step 1: Define the Counterfactual Probabilities**} \\
\texttt{We need to evaluate \$P(M\_\{uvq\} = 1 | Y\_\{upt\}=1, Z\_\{upj\} = z)\$ and \$P(M\_\{uvq\} = 1 | Y\_\{upt\}=0, Z\_\{upj\} = z)\$ for the relevant values of \$z\$ (\$0\$ and \$1\$).} \\
\\ 
\texttt{From the problem statement, we have the following conditional probabilities for \$M\_\{uvq\}\$:} \\
\texttt{*   \$P(M=1 | Y=0, Z=0) = 44\textbackslash{}\%\$} \\
\texttt{*   \$P(M=1 | Y=0, Z=1) = 51\textbackslash{}\%\$} \\
\texttt{*   \$P(M=1 | Y=1, Z=0) = 57\textbackslash{}\%\$} \\
\texttt{*   \$P(M=1 | Y=1, Z=1) = 10\textbackslash{}\%\$} \\
\\ 
\texttt{**Step 2: Calculate the Indirect Effect (NIE)**} \\
\texttt{The indirect effect measures the change in \$M\$ due to the change in \$Z\$ caused by changing \$Y\$, while keeping \$Y\$ fixed at the baseline (0).} \\
\texttt{*   Effect of \$Z\$ on \$M\$ when \$Y=0\$:} \\
\texttt{    *   If \$Z\$ goes from 0 to 1 (while \$Y=0\$): \$P(M=1|Y=0, Z=1) {-} P(M=1|Y=0, Z=0) = 51\textbackslash{}\% {-} 44\textbackslash{}\% = +7\textbackslash{}\%\$.} \\
\texttt{    *   This indicates a **positive** association between \$Z\$ and \$M\$ in the absence of \$Y\$.} \\
\\ 
\texttt{**Step 3: Calculate the Direct Effect (NDE)**} \\
\texttt{The direct effect measures the change} \\
            }
        }
    }
    & \\ \\

    \theutterance \stepcounter{utterance}  
    & & \multicolumn{4}{p{0.6\linewidth}}{
        \cellcolor[rgb]{0.95,0.95,0.95}{
            \makecell[{{p{\linewidth}}}]{
                \texttt{\tiny{[P1$\leftrightarrow$P]}}
                \texttt{{[}SELF{-}FEEDBACK{]}} \\
\texttt{FORMAT: problem} \\
\texttt{RULES: problem} \\
\texttt{STRATEGY: problem} \\
\texttt{REASON: The answer fails the format requirement by not starting with 'Yes' or 'No', incorrectly identifies the causal graph structure by omitting the direct effect of Yupt on Muvq, and leaves the reasoning incomplete.} \\
            }
        }
    }
    & \\ \\

    \theutterance \stepcounter{utterance}  
    & \multicolumn{4}{p{0.6\linewidth}}{
        \cellcolor[rgb]{0.561,0.737,0.561}{
            \makecell[{{p{\linewidth}}}]{
                \texttt{\tiny{[P1$\rangle$GM]}}
                \texttt{No} \\
\\ 
\texttt{**Reasoning:**} \\
\\ 
\texttt{To determine if \$Y\_\{upt\}\$ negatively affects \$M\_\{uvq\}\$ through \$Z\_\{upj\}\$, we must analyze the **Natural Indirect Effect (NIE)**, which isolates the effect of changing \$Y\_\{upt\}\$ on \$M\_\{uvq\}\$ mediated solely by \$Z\_\{upj\}\$. This is calculated as the difference in the probability of \$M\_\{uvq\}\$ when \$Z\_\{upj\}\$ changes from its baseline value (under \$Y\_\{upt\}=0\$) to its intervened value (under \$Y\_\{upt\}=1\$), while holding \$Y\_\{upt\}\$ fixed at 0.} \\
\\ 
\texttt{The causal graph structure is:} \\
\texttt{\$X\_\{yfo\} \textbackslash{}to Z\_\{upj\} \textbackslash{}to M\_\{uvq\}\$} \\
\texttt{\$Y\_\{upt\} \textbackslash{}to Z\_\{upj\}\$} \\
\texttt{\$Y\_\{upt\} \textbackslash{}to M\_\{uvq\}\$} \\
\\ 
\texttt{**Step 1: Identify Conditional Probabilities for \$M\_\{uvq\}\$**} \\
\texttt{From the problem statement:} \\
\texttt{*   \$P(M=1 | Y=0, Z=0) = 44\textbackslash{}\%\$} \\
\texttt{*   \$P(M=1 | Y=0, Z=1) = 51\textbackslash{}\%\$} \\
\texttt{*   \$P(M=1 | Y=1, Z=0) = 57\textbackslash{}\%\$} \\
\texttt{*   \$P(M=1 | Y=1, Z=1) = 10\textbackslash{}\%\$} \\
\\ 
\texttt{**Step 2: Calculate the Natural Indirect Effect (NIE)**} \\
\texttt{The NIE is defined as \$E{[}M(1, Z(0)) {-} M(1, Z(1)){]}\$ or equivalently \$P(M=1 | Y=1, Z=0) {-} P(M=1 | Y=1, Z=1)\$. This measures the change in \$M\$ when \$Z\$ switches from the state it would take if \$Y=0\$ (which is \$Z=0\$ in the counterfactual logic for the specific comparison of the mediator's effect) to the state it takes when \$Y=1\$, holding \$Y\$ constant at 1.} \\
\\ 
\texttt{However, a more direct interpretation of "Does Y affect M through Z" requires comparing the outcome when the mediator is forced to the value it would have taken under \$Y=0\$ versus the value it takes under \$Y=1\$, while \$Y\$ is set to 1.} \\
\texttt{*   Probability of \$M\$ when \$Y=1\$ and \$Z\$ is set to the value it would have been if \$Y=0\$ (which is \$Z=0\$ for the baseline group, but we need the distribution of \$Z\$).} \\
\texttt{*   Actually, the standard definition for the NIE in a binary setting is: \$P(M=1 | Y=1, Z=} \\
            }
        }
    }
    & & \\ \\

    \theutterance \stepcounter{utterance}  
    & & & \multicolumn{2}{p{0.3\linewidth}}{
        \cellcolor[rgb]{0.9,0.9,0.9}{
            \makecell[{{p{\linewidth}}}]{
                \texttt{\tiny{[GM$|$GM]}}
                \texttt{No} \\
            }
        }
    }
    & & \\ \\

    \theutterance \stepcounter{utterance}  
    & & & \multicolumn{2}{p{0.3\linewidth}}{
        \cellcolor[rgb]{0.9,0.9,0.9}{
            \makecell[{{p{\linewidth}}}]{
                \texttt{\tiny{[GM$|$GM]}}
                \texttt{yes} \\
            }
        }
    }
    & & \\ \\

    \theutterance \stepcounter{utterance}  
    & & & \multicolumn{2}{p{0.3\linewidth}}{
        \cellcolor[rgb]{0.9,0.9,0.9}{
            \makecell[{{p{\linewidth}}}]{
                \texttt{\tiny{[GM$|$GM]}}
                \texttt{game\_result = LOSE} \\
            }
        }
    }
    & & \\ \\

\end{supertabular}
}

\end{document}
